\documentclass[12pt,a4paper]{article}

% ----------------------------------------------------------------
% Кодировка, язык, шрифты
% ----------------------------------------------------------------
\usepackage[T2A]{fontenc}
\usepackage[utf8]{inputenc}
\usepackage[russian]{babel}

% ----------------------------------------------------------------
% Математика
% ----------------------------------------------------------------
\usepackage{amsmath,amssymb,mathtools,amsthm}

% ----------------------------------------------------------------
% Оформление
% ----------------------------------------------------------------
\usepackage{geometry}
\usepackage{booktabs}
\usepackage{array}
\usepackage{tabularx}
\usepackage{longtable}
\usepackage{enumitem}
\usepackage{xcolor}
\usepackage{tikz}
\usetikzlibrary{arrows.meta,positioning,calc}

\geometry{
    left=2cm,
    right=2cm,
    top=2cm,
    bottom=2cm
}

\setlength{\parindent}{1.25em}
\setlength{\parskip}{0.25em}
\setlist[itemize]{topsep=3pt,itemsep=2pt}
\setlist[enumerate]{topsep=3pt,itemsep=2pt}

% ----------------------------------------------------------------
% Теоремы и определения
% ----------------------------------------------------------------
\newtheorem{theorem}{Теорема}[section]
\newtheorem{lemma}[theorem]{Лемма}
\newtheorem{proposition}[theorem]{Утверждение}
\newtheorem{corollary}[theorem]{Следствие}

\theoremstyle{definition}
\newtheorem{definition}[theorem]{Определение}
\newtheorem{example}[theorem]{Пример}

\theoremstyle{remark}
\newtheorem{remark}[theorem]{Замечание}

% ----------------------------------------------------------------
% Обозначения
% ----------------------------------------------------------------
\newcommand{\R}{\mathbb{R}}
\newcommand{\E}{\mathbb{E}}
\newcommand{\PP}{\mathbb{P}}
\newcommand{\Var}{\operatorname{Var}}
\newcommand{\Cov}{\operatorname{Cov}}
\newcommand{\rank}{\operatorname{rank}}
\newcommand{\diag}{\operatorname{diag}}
\newcommand{\tr}{\operatorname{tr}}
\newcommand{\col}{\operatorname{col}}
\newcommand{\sign}{\operatorname{sign}}
\newcommand{\Ind}{\mathbf{1}}

\DeclareMathOperator*{\argmin}{arg\,min}
\DeclareMathOperator*{\argmax}{arg\,max}

% ----------------------------------------------------------------
% Блоки алгоритмов
% ----------------------------------------------------------------
\newenvironment{algorithmbox}[1]
{
    \begin{center}
    \begin{minipage}{0.96\textwidth}
    \hrule
    \medskip
    \textbf{Алгоритм: #1.}
    \par\smallskip
}
{
    \medskip
    \hrule
    \end{minipage}
    \end{center}
}

% ----------------------------------------------------------------

\title{Билет №94 \\ \small Методы машинного обучения. Линейная регрессия. Кластеризация. (ИТМО)}
\author{Сериков Владислав Александрович}
\date{19 июля 2026}

\begin{document}

\maketitle
\tableofcontents
\newpage

% ================================================================
\section{Основные понятия машинного обучения}
% ================================================================

\subsection{Постановка задачи}

\begin{definition}
\emph{Машинное обучение} --- совокупность методов построения моделей,
которые извлекают закономерности из данных и используют их для предсказания,
описания структуры данных или принятия решений.
\end{definition}

Один объект описывается вектором признаков
\[
    x=(x_1,\ldots,x_d)^\top\in\mathcal{X}.
\]
В задачах обучения с учителем объекту дополнительно соответствует целевая
переменная
\[
    y\in\mathcal{Y}.
\]
Набор наблюдений имеет вид
\[
    D=\{(x_i,y_i)\}_{i=1}^{n}.
\]

Модель представляет собой параметризованную функцию
\[
    f_\theta\colon \mathcal{X}\to\mathcal{Y},
\]
где $\theta$ --- параметры, определяемые по обучающим данным.

Следует различать:

\begin{itemize}
    \item \textbf{параметры модели} $\theta$, которые оцениваются при обучении;
    \item \textbf{гиперпараметры}, задающие структуру или режим обучения модели:
          коэффициент регуляризации, число кластеров, глубину дерева,
          скорость обучения и т.\,д.;
    \item \textbf{признаки} --- числовые или категориальные характеристики объектов;
    \item \textbf{целевую переменную} --- величину, которую требуется предсказать.
\end{itemize}

\subsection{Основные типы обучения}

\begin{table}[ht]
\centering
\begin{tabularx}{\textwidth}{
    >{\raggedright\arraybackslash}p{3.2cm}
    >{\raggedright\arraybackslash}X
    >{\raggedright\arraybackslash}p{4.4cm}}
\toprule
\textbf{Тип обучения} & \textbf{Постановка} & \textbf{Типичные задачи}\\
\midrule
Обучение с учителем &
Даны пары $(x_i,y_i)$. Требуется восстановить зависимость $y$ от $x$. &
Регрессия, бинарная и многоклассовая классификация.\\

Обучение без учителя &
Даны только объекты $x_i$. Требуется обнаружить скрытую структуру данных. &
Кластеризация, снижение размерности, оценивание плотности, поиск аномалий.\\

Частично контролируемое обучение &
Имеется небольшое количество размеченных и большое количество
неразмеченных объектов. &
Классификация при высокой стоимости получения разметки.\\

Самообучение &
Целевой сигнал автоматически конструируется из самих данных. &
Предобучение языковых, визуальных и мультимодальных моделей.\\

Обучение с подкреплением &
Агент взаимодействует со средой, выбирает действия и получает награды. &
Управление, робототехника, игры, последовательное принятие решений.\\
\bottomrule
\end{tabularx}
\end{table}

В обучении с учителем выделяют две основные задачи.

\begin{definition}
\emph{Регрессия} --- задача предсказания числовой целевой переменной
$y\in\R$.
\end{definition}

\begin{definition}
\emph{Классификация} --- задача предсказания конечного класса
$y\in\{1,\ldots,K\}$ либо вероятностей принадлежности классам.
\end{definition}

\subsection{Основные семейства моделей}

\begin{itemize}
    \item \textbf{Линейные модели} представляют прогноз как линейную комбинацию
          признаков. Они вычислительно эффективны и хорошо интерпретируются.
    \item \textbf{Модели на основе расстояний} сравнивают новый объект
          с похожими объектами обучающей выборки.
    \item \textbf{Деревья решений} последовательно разбивают пространство
          признаков с помощью логических условий.
    \item \textbf{Ансамбли} объединяют несколько моделей. Основные идеи ---
          усреднение, построение моделей на случайных подвыборках и
          последовательное исправление ошибок.
    \item \textbf{Ядерные методы} неявно отображают объекты в пространство
          большей размерности с помощью функции сходства --- ядра.
    \item \textbf{Нейронные сети} строят композиции линейных преобразований
          и нелинейных функций активации.
    \item \textbf{Вероятностные модели} задают распределение данных и позволяют
          оценивать не только прогноз, но и связанную с ним неопределённость.
\end{itemize}

Дополнительно различают:

\begin{itemize}
    \item \textbf{параметрические методы}, в которых число параметров заранее
          фиксировано;
    \item \textbf{непараметрические методы}, сложность которых может возрастать
          вместе с объёмом данных;
    \item \textbf{дискриминативные модели}, непосредственно моделирующие
          зависимость $y$ от $x$;
    \item \textbf{генеративные модели}, описывающие распределение данных
          $p(x)$ или совместное распределение $p(x,y)$.
\end{itemize}

\subsection{Функция потерь и эмпирический риск}

\begin{definition}
\emph{Функция потерь} $L(y,\widehat y)$ измеряет ошибку прогноза
$\widehat y$ относительно истинного значения $y$.
\end{definition}

Примеры функций потерь:

\begin{align*}
    L_{\mathrm{MSE}}(y,\widehat y)
        &= (y-\widehat y)^2,\\
    L_{\mathrm{MAE}}(y,\widehat y)
        &= |y-\widehat y|,\\
    L_{\mathrm{log}}(y,\widehat p)
        &= -y\log\widehat p-(1-y)\log(1-\widehat p).
\end{align*}

Если $(X,Y)$ --- случайная пара из генеральной совокупности, то
\emph{теоретический риск} равен
\[
    R(f)=\E L(Y,f(X)).
\]
Распределение данных обычно неизвестно, поэтому используется
\emph{эмпирический риск}
\[
    \widehat R_n(f)
    =
    \frac{1}{n}\sum_{i=1}^{n}L(y_i,f(x_i)).
\]

Принцип минимизации эмпирического риска имеет вид
\[
    \widehat f
    \in
    \argmin_{f\in\mathcal{F}}\widehat R_n(f),
\]
где $\mathcal{F}$ --- выбранный класс моделей.

Чтобы ограничить сложность модели, вводится регуляризация:
\[
    \widehat\theta
    \in
    \argmin_{\theta}
    \left\{
        \frac{1}{n}\sum_{i=1}^{n}
        L(y_i,f_\theta(x_i))
        +
        \lambda\Omega(\theta)
    \right\}.
\]
Здесь $\Omega(\theta)$ --- штраф за сложность, а $\lambda\geq 0$ ---
коэффициент регуляризации.

\subsection{Недообучение и переобучение}

\begin{definition}
\emph{Недообучение} возникает, когда класс моделей слишком прост и модель
не способна описать существенные закономерности даже в обучающей выборке.
\end{definition}

\begin{definition}
\emph{Переобучение} возникает, когда модель хорошо описывает обучающие данные,
включая случайный шум, но показывает низкое качество на новых объектах.
\end{definition}

Типичные способы борьбы с переобучением:

\begin{itemize}
    \item увеличение объёма и разнообразия данных;
    \item уменьшение сложности модели;
    \item регуляризация;
    \item отбор признаков;
    \item ранняя остановка;
    \item перекрёстная проверка;
    \item ансамблирование;
    \item контроль утечки информации.
\end{itemize}

\subsection{Разложение ошибки на смещение и дисперсию}

\begin{theorem}[Разложение квадратичной ошибки]
Пусть для фиксированного объекта $x$ выполняется
\[
    Y=f_*(x)+\varepsilon,
    \qquad
    \E\varepsilon=0,
    \qquad
    \Var(\varepsilon)=\sigma^2.
\]
Пусть $\widehat f_D(x)$ --- прогноз модели, построенной по случайной
обучающей выборке $D$, причём $D$ не зависит от ошибки $\varepsilon$.
Тогда
\[
\begin{split}
    \E_{D,\varepsilon}
    \left[
        \bigl(Y-\widehat f_D(x)\bigr)^2
    \right]
    &=
    \sigma^2
    +
    \left(
        \E_D\widehat f_D(x)-f_*(x)
    \right)^2\\
    &\quad+
    \Var_D\bigl(\widehat f_D(x)\bigr).
\end{split}
\]
\end{theorem}

\begin{proof}
Обозначим
\[
    m(x)=\E_D\widehat f_D(x).
\]
Тогда
\[
\begin{split}
    Y-\widehat f_D(x)
    &=
    f_*(x)+\varepsilon-\widehat f_D(x)\\
    &=
    \bigl(f_*(x)-m(x)\bigr)
    +
    \bigl(m(x)-\widehat f_D(x)\bigr)
    +
    \varepsilon.
\end{split}
\]
Возведём выражение в квадрат и возьмём математическое ожидание. Смешанные
слагаемые равны нулю, поскольку
\[
    \E_D\bigl[m(x)-\widehat f_D(x)\bigr]=0,
    \qquad
    \E\varepsilon=0,
\]
а $\varepsilon$ не зависит от обучающей выборки. Следовательно,
\[
\begin{split}
    \E_{D,\varepsilon}
    \left[
        \bigl(Y-\widehat f_D(x)\bigr)^2
    \right]
    &=
    \bigl(f_*(x)-m(x)\bigr)^2\\
    &\quad+
    \E_D
    \left[
        \bigl(\widehat f_D(x)-m(x)\bigr)^2
    \right]
    +
    \E\varepsilon^2.
\end{split}
\]
Второе слагаемое является дисперсией прогноза, а
$\E\varepsilon^2=\sigma^2$. Получаем требуемое разложение.
\end{proof}

Таким образом, средняя ошибка состоит из:

\begin{itemize}
    \item неустранимого шума $\sigma^2$;
    \item квадрата смещения;
    \item дисперсии модели.
\end{itemize}

Слишком простые модели обычно имеют большое смещение, а слишком сложные ---
большую дисперсию.

\subsection{Разбиение данных и проверка качества}

Обычно данные разделяют на:

\begin{itemize}
    \item \textbf{обучающую выборку} --- для оценки параметров;
    \item \textbf{валидационную выборку} --- для выбора гиперпараметров;
    \item \textbf{тестовую выборку} --- для окончательной независимой оценки.
\end{itemize}

\begin{figure}[ht]
\centering
\begin{tikzpicture}[
    >=Latex,
    block/.style={
        draw,
        rounded corners,
        align=center,
        minimum height=9mm,
        text width=2.6cm
    },
    node distance=5mm
]
\node[block] (data) {Исходные\\данные};
\node[block, right=of data] (split) {Разбиение\\выборки};
\node[block, right=of split] (train) {Обучение\\модели};
\node[block, right=of train] (valid) {Выбор\\гиперпараметров};
\node[block, right=of valid] (test) {Оценка на\\тесте};

\draw[->] (data) -- (split);
\draw[->] (split) -- (train);
\draw[->] (train) -- (valid);
\draw[->] (valid) -- (test);
\end{tikzpicture}
\caption{Общая схема построения модели машинного обучения.}
\end{figure}

\begin{definition}
\emph{Утечка данных} --- попадание в процесс обучения информации, которая
не должна быть доступна модели в момент применения.
\end{definition}

Например, стандартизацию признаков необходимо настраивать только по
обучающей части данных. Использование среднего и дисперсии всей выборки
приводит к утечке информации из валидационной или тестовой части.

\begin{algorithmbox}{перекрёстная проверка по $K$ блокам}
\begin{enumerate}
    \item Разделить обучающую выборку на $K$ непересекающихся блоков
          $D_1,\ldots,D_K$.
    \item Для каждого $j=1,\ldots,K$:
    \begin{enumerate}
        \item обучить модель на объединении всех блоков, кроме $D_j$;
        \item вычислить метрику на блоке $D_j$.
    \end{enumerate}
    \item Усреднить полученные значения метрики.
    \item Выбрать гиперпараметры по среднему качеству.
    \item Заново обучить модель с выбранными гиперпараметрами на всей
          обучающей выборке.
    \item Один раз оценить итоговое качество на тестовой выборке.
\end{enumerate}
\end{algorithmbox}

\begin{example}
Пусть имеется шесть объектов и $K=3$. Можно взять блоки
\[
    D_1=\{1,2\},\qquad
    D_2=\{3,4\},\qquad
    D_3=\{5,6\}.
\]
Каждый объект один раз используется для проверки и два раза попадает
в обучающую часть.
\end{example}

Для временных рядов нельзя случайно перемешивать прошлые и будущие
наблюдения. Для зависимых групп объектов следует разбивать данные так,
чтобы связанные объекты не попадали одновременно в обучение и проверку.

% ================================================================
\section{Линейная регрессия}
% ================================================================

\subsection{Модель линейной регрессии}

Пусть объект имеет $d$ исходных признаков. После добавления единичного
признака запишем
\[
    \widetilde x_i
    =
    \begin{pmatrix}
        1\\x_{i1}\\ \vdots\\x_{id}
    \end{pmatrix}
    \in\R^p,
    \qquad p=d+1.
\]
Модель линейной регрессии задаётся формулой
\[
    \widehat y_i
    =
    \widetilde x_i^\top\beta
    =
    \beta_0+\beta_1x_{i1}+\cdots+\beta_dx_{id}.
\]

Параметр $\beta_0$ называется свободным членом, а $\beta_j$ при $j\geq 1$ ---
коэффициентом соответствующего признака.

Соберём объекты в матрицу признаков:
\[
    X=
    \begin{pmatrix}
        \widetilde x_1^\top\\
        \vdots\\
        \widetilde x_n^\top
    \end{pmatrix}
    \in\R^{n\times p},
    \qquad
    y=
    \begin{pmatrix}
        y_1\\ \vdots\\y_n
    \end{pmatrix}.
\]
Тогда вектор прогнозов равен
\[
    \widehat y=X\beta.
\]

Термин ``линейная'' означает линейность по параметрам. Например,
\[
    y
    =
    \beta_0+\beta_1x+\beta_2x^2+\varepsilon
\]
также является линейной регрессией относительно параметров
$\beta_0,\beta_1,\beta_2$.

\subsection{Метод наименьших квадратов}

Остаток на $i$-м объекте:
\[
    e_i=y_i-\widehat y_i.
\]
Вектор остатков:
\[
    e=y-X\beta.
\]

Метод наименьших квадратов выбирает параметры, минимизирующие сумму
квадратов остатков:
\[
    \widehat\beta
    \in
    \argmin_{\beta\in\R^p}
    \operatorname{RSS}(\beta),
    \qquad
    \operatorname{RSS}(\beta)
    =
    \|y-X\beta\|_2^2.
\]

Градиент и матрица Гессе имеют вид
\[
    \nabla_\beta\operatorname{RSS}(\beta)
    =
    2X^\top(X\beta-y),
\]
\[
    \nabla_\beta^2\operatorname{RSS}(\beta)
    =
    2X^\top X.
\]
Матрица $X^\top X$ неотрицательно определена, поэтому задача является
выпуклой.

\begin{theorem}[Нормальные уравнения]
Для любых $X\in\R^{n\times p}$ и $y\in\R^n$ вектор
$\widehat\beta$ минимизирует
\[
    \|y-X\beta\|_2^2
\]
тогда и только тогда, когда он удовлетворяет нормальным уравнениям
\[
    X^\top X\widehat\beta=X^\top y.
\]
Минимум всегда существует. Если $\rank X=p$, решение единственно и равно
\[
    \boxed{
    \widehat\beta=(X^\top X)^{-1}X^\top y
    }.
\]
В общем случае множество решений можно записать через псевдообратную
матрицу Мура--Пенроуза:
\[
    \widehat\beta
    =
    X^+y+(I-X^+X)z,
    \qquad z\in\R^p.
\]
\end{theorem}

\begin{proof}
Функция
\[
    F(\beta)=\|y-X\beta\|_2^2
\]
дифференцируема и выпукла. Поэтому точка минимума должна удовлетворять
условию
\[
    \nabla F(\widehat\beta)=0.
\]
Так как
\[
    \nabla F(\beta)=2X^\top(X\beta-y),
\]
получаем
\[
    X^\top X\widehat\beta=X^\top y.
\]

Докажем достаточность. Пусть $\widehat\beta$ удовлетворяет нормальным
уравнениям, и обозначим
\[
    r=y-X\widehat\beta.
\]
Тогда
\[
    X^\top r=0.
\]
Для произвольного $h\in\R^p$ имеем
\[
\begin{split}
    \|y-X(\widehat\beta+h)\|_2^2
    &=
    \|r-Xh\|_2^2\\
    &=
    \|r\|_2^2
    -2r^\top Xh
    +
    \|Xh\|_2^2.
\end{split}
\]
Но
\[
    r^\top Xh=(X^\top r)^\top h=0.
\]
Следовательно,
\[
    \|y-X(\widehat\beta+h)\|_2^2
    =
    \|r\|_2^2+\|Xh\|_2^2
    \geq
    \|r\|_2^2.
\]
Значит, $\widehat\beta$ является точкой глобального минимума.

Существование следует из существования ортогональной проекции
$y$ на конечномерное подпространство $\col(X)$. Пусть эта проекция равна
$\widehat y\in\col(X)$. Тогда существует $\widehat\beta$, для которого
$X\widehat\beta=\widehat y$, и остаток $y-\widehat y$ ортогонален
$\col(X)$. Поэтому
\[
    X^\top(y-X\widehat\beta)=0.
\]

Если $\rank X=p$, матрица $X^\top X$ положительно определена:
для любого $v\ne 0$
\[
    v^\top X^\top Xv=\|Xv\|_2^2>0.
\]
Следовательно, она обратима, и решение единственно.

Если матрица $X$ имеет неполный столбцовый ранг, все минимизаторы дают
один и тот же вектор прогнозов, а их разности принадлежат ядру $X$.
Стандартная параметризация всех решений имеет вид
\[
    X^+y+(I-X^+X)z.
\]
\end{proof}

\subsection{Геометрический смысл метода наименьших квадратов}

Пусть $\rank X=p$. Определим
\[
    P_X=X(X^\top X)^{-1}X^\top.
\]

\begin{proposition}
Матрица $P_X$ является матрицей ортогонального проектирования на
столбцовое пространство $\col(X)$. В частности,
\[
    P_X^\top=P_X,
    \qquad
    P_X^2=P_X.
\]
Вектор прогнозов и остатки равны
\[
    \widehat y=P_Xy,
    \qquad
    e=(I-P_X)y,
\]
причём
\[
    X^\top e=0.
\]
\end{proposition}

\begin{proof}
Симметричность следует из равенства
\[
\begin{split}
    P_X^\top
    &=
    \left(
        X(X^\top X)^{-1}X^\top
    \right)^\top\\
    &=
    X(X^\top X)^{-1}X^\top=P_X,
\end{split}
\]
поскольку $X^\top X$ и его обратная матрица симметричны.

Далее,
\[
\begin{split}
    P_X^2
    &=
    X(X^\top X)^{-1}X^\top
    X(X^\top X)^{-1}X^\top\\
    &=
    X(X^\top X)^{-1}X^\top=P_X.
\end{split}
\]
Следовательно, $P_X$ является симметричной идемпотентной матрицей, то есть
ортогональным проектором.

Из формулы для оценки параметров получаем
\[
    \widehat y
    =
    X\widehat\beta
    =
    X(X^\top X)^{-1}X^\top y
    =
    P_Xy.
\]
Наконец,
\[
    X^\top e
    =
    X^\top(y-P_Xy)
    =
    X^\top y-X^\top X(X^\top X)^{-1}X^\top y
    =
    0.
\]
\end{proof}

Таким образом, метод наименьших квадратов ищет ближайший к $y$ вектор,
лежащий в линейном пространстве $\col(X)$.

\subsection{Минимальный численный пример}

Рассмотрим наблюдения
\[
    (x_1,y_1)=(1,2),\qquad
    (x_2,y_2)=(2,3),\qquad
    (x_3,y_3)=(3,5).
\]
Матрица признаков со свободным членом и вектор ответов:
\[
    X=
    \begin{pmatrix}
        1&1\\
        1&2\\
        1&3
    \end{pmatrix},
    \qquad
    y=
    \begin{pmatrix}
        2\\3\\5
    \end{pmatrix}.
\]
Вычислим
\[
    X^\top X
    =
    \begin{pmatrix}
        3&6\\
        6&14
    \end{pmatrix},
    \qquad
    X^\top y
    =
    \begin{pmatrix}
        10\\23
    \end{pmatrix}.
\]
Так как
\[
    (X^\top X)^{-1}
    =
    \frac{1}{6}
    \begin{pmatrix}
        14&-6\\
        -6&3
    \end{pmatrix},
\]
получаем
\[
    \widehat\beta
    =
    \begin{pmatrix}
        1/3\\3/2
    \end{pmatrix}.
\]
Искомая прямая:
\[
    \boxed{
    \widehat y=\frac{1}{3}+\frac{3}{2}x
    }.
\]

Прогнозы и остатки:
\[
    \widehat y=
    \begin{pmatrix}
        11/6\\10/3\\29/6
    \end{pmatrix},
    \qquad
    e=
    \begin{pmatrix}
        1/6\\-1/3\\1/6
    \end{pmatrix}.
\]
Сумма квадратов остатков:
\[
    \operatorname{RSS}
    =
    \frac{1}{36}+\frac{1}{9}+\frac{1}{36}
    =
    \frac{1}{6}.
\]

\begin{figure}[ht]
\centering
\begin{tikzpicture}[scale=0.9]
    \draw[->] (0,0) -- (4.3,0) node[right] {$x$};
    \draw[->] (0,0) -- (0,6.3) node[above] {$y$};

    \foreach \x in {1,2,3,4}
        \draw (\x,0.08)--(\x,-0.08) node[below] {\x};

    \foreach \y in {1,2,3,4,5,6}
        \draw (0.08,\y)--(-0.08,\y) node[left] {\y};

    \fill (1,2) circle (2.5pt);
    \fill (2,3) circle (2.5pt);
    \fill (3,5) circle (2.5pt);

    \draw[blue,thick] (0,0.333) -- (3.8,6.033);
    \node[blue,right] at (2.5,4.2)
        {$\widehat y=\frac13+\frac32x$};
\end{tikzpicture}
\caption{Линейная аппроксимация трёх наблюдений.}
\end{figure}

\subsection{Численное решение}

Явное вычисление матрицы $(X^\top X)^{-1}$ удобно для теоретического
анализа, но может быть численно неустойчивым. Обусловленность матрицы
$X^\top X$ примерно равна квадрату обусловленности $X$.

\begin{algorithmbox}{устойчивое обучение линейной регрессии}
\begin{enumerate}
    \item Сформировать матрицу признаков $X$ и добавить столбец единиц.
    \item При необходимости закодировать категориальные признаки и
          стандартизовать числовые признаки.
    \item Решить задачу наименьших квадратов:
    \begin{itemize}
        \item с помощью QR-разложения $X=QR$ и системы
              \[
                  R\widehat\beta=Q^\top y;
              \]
        \item либо с помощью сингулярного разложения при вырожденности
              или плохой обусловленности.
    \end{itemize}
    \item Вычислить прогнозы $\widehat y=X\widehat\beta$.
    \item Вычислить остатки и метрики качества.
    \item Проверить модель на независимых данных.
\end{enumerate}
\end{algorithmbox}

\subsection{Градиентный спуск}

Рассмотрим нормированную целевую функцию
\[
    J(\beta)
    =
    \frac{1}{2n}\|X\beta-y\|_2^2.
\]
Её градиент равен
\[
    \nabla J(\beta)
    =
    \frac{1}{n}X^\top(X\beta-y).
\]

\begin{algorithmbox}{градиентный спуск для линейной регрессии}
\begin{enumerate}
    \item Выбрать начальное значение $\beta^{(0)}$ и скорость обучения
          $\eta>0$.
    \item Для $t=0,1,2,\ldots$ вычислять
          \[
              g^{(t)}
              =
              \frac{1}{n}X^\top
              \bigl(X\beta^{(t)}-y\bigr).
          \]
    \item Обновлять параметры:
          \[
              \beta^{(t+1)}
              =
              \beta^{(t)}-\eta g^{(t)}.
          \]
    \item Остановиться, если изменение целевой функции или параметров стало
          меньше заданного порога.
\end{enumerate}
\end{algorithmbox}

\begin{theorem}[Сходимость градиентного спуска]
Пусть $\rank X=p$ и
\[
    H=\frac{1}{n}X^\top X.
\]
Обозначим через $\lambda_{\max}(H)$ наибольшее собственное значение $H$.
Если
\[
    0<\eta<\frac{2}{\lambda_{\max}(H)},
\]
то итерации градиентного спуска сходятся к единственному решению метода
наименьших квадратов.
\end{theorem}

\begin{proof}
Пусть $\widehat\beta$ --- решение нормальных уравнений. Тогда
\[
    H\widehat\beta=\frac{1}{n}X^\top y.
\]
Обозначим ошибку параметров
\[
    e^{(t)}=\beta^{(t)}-\widehat\beta.
\]
Итерация градиентного спуска имеет вид
\[
\begin{split}
    \beta^{(t+1)}
    &=
    \beta^{(t)}
    -
    \eta
    \left(
        H\beta^{(t)}-\frac{1}{n}X^\top y
    \right)\\
    &=
    \beta^{(t)}-\eta H
    \bigl(\beta^{(t)}-\widehat\beta\bigr).
\end{split}
\]
Следовательно,
\[
    e^{(t+1)}=(I-\eta H)e^{(t)}.
\]

Матрица $H$ симметрична и положительно определена, поэтому существует
ортогональная матрица $U$ и положительные собственные значения
$\lambda_1,\ldots,\lambda_p$ такие, что
\[
    H=U\diag(\lambda_1,\ldots,\lambda_p)U^\top.
\]
Переходя к координатам $z^{(t)}=U^\top e^{(t)}$, получаем
\[
    z_j^{(t+1)}
    =
    (1-\eta\lambda_j)z_j^{(t)}.
\]
Условие
\[
    0<\eta<\frac{2}{\lambda_{\max}(H)}
\]
гарантирует
\[
    |1-\eta\lambda_j|<1
\]
для всех $j$. Поэтому каждая координата $z_j^{(t)}$ стремится к нулю,
а значит,
\[
    \beta^{(t)}\to\widehat\beta.
\]
\end{proof}

\begin{example}
Пусть
\[
    J(\beta)=\frac12(\beta-2)^2,
    \qquad
    \eta=\frac12,
    \qquad
    \beta^{(0)}=0.
\]
Тогда
\[
    \beta^{(t+1)}
    =
    \beta^{(t)}-\frac12(\beta^{(t)}-2).
\]
Получаем
\[
    \beta^{(1)}=1,\qquad
    \beta^{(2)}=1.5,\qquad
    \beta^{(3)}=1.75.
\]
Последовательность сходится к точному решению $\widehat\beta=2$.
\end{example}

\subsection{Статистическая модель}

Классическая линейная модель записывается как
\[
    y=X\beta+\varepsilon,
\]
где
\[
    \varepsilon=
    \begin{pmatrix}
        \varepsilon_1\\ \vdots\\\varepsilon_n
    \end{pmatrix}
\]
--- случайный вектор ошибок.

Основные предположения:

\begin{enumerate}
    \item \textbf{Линейность по параметрам:}
          \[
              \E[y\mid X]=X\beta.
          \]
    \item \textbf{Экзогенность:}
          \[
              \E[\varepsilon\mid X]=0.
          \]
    \item \textbf{Отсутствие точной мультиколлинеарности:}
          \[
              \rank X=p.
          \]
    \item \textbf{Гомоскедастичность и отсутствие корреляции ошибок:}
          \[
              \Cov(\varepsilon\mid X)=\sigma^2I_n.
          \]
    \item \textbf{Нормальность} для точного конечновыборочного статистического
          вывода:
          \[
              \varepsilon\mid X\sim
              \mathcal{N}(0,\sigma^2I_n).
          \]
\end{enumerate}

Нормальность не требуется для вычисления оценки и для теоремы
Гаусса--Маркова.

\begin{proposition}
При условиях
\[
    y=X\beta+\varepsilon,
    \qquad
    \E[\varepsilon\mid X]=0,
    \qquad
    \Cov(\varepsilon\mid X)=\sigma^2I_n
\]
оценка метода наименьших квадратов несмещённа:
\[
    \E[\widehat\beta\mid X]=\beta,
\]
а её условная ковариационная матрица равна
\[
    \Cov(\widehat\beta\mid X)
    =
    \sigma^2(X^\top X)^{-1}.
\]
\end{proposition}

\begin{proof}
Имеем
\[
\begin{split}
    \widehat\beta
    &=
    (X^\top X)^{-1}X^\top y\\
    &=
    (X^\top X)^{-1}X^\top(X\beta+\varepsilon)\\
    &=
    \beta+(X^\top X)^{-1}X^\top\varepsilon.
\end{split}
\]
Поэтому
\[
    \E[\widehat\beta\mid X]
    =
    \beta+
    (X^\top X)^{-1}X^\top
    \E[\varepsilon\mid X]
    =
    \beta.
\]

Для ковариации получаем
\[
\begin{split}
    \Cov(\widehat\beta\mid X)
    &=
    (X^\top X)^{-1}X^\top
    \Cov(\varepsilon\mid X)
    X(X^\top X)^{-1}\\
    &=
    \sigma^2
    (X^\top X)^{-1}X^\top X(X^\top X)^{-1}\\
    &=
    \sigma^2(X^\top X)^{-1}.
\end{split}
\]
\end{proof}

\subsection{Связь метода наименьших квадратов с максимальным правдоподобием}

\begin{theorem}
Пусть наблюдения условно независимы и
\[
    y_i\mid x_i
    \sim
    \mathcal{N}(x_i^\top\beta,\sigma^2),
    \qquad
    \sigma^2>0.
\]
При фиксированном $\sigma^2$ оценки максимального правдоподобия параметра
$\beta$ совпадают с оценками метода наименьших квадратов.
\end{theorem}

\begin{proof}
Функция правдоподобия равна
\[
    L(\beta)
    =
    \prod_{i=1}^{n}
    \frac{1}{\sqrt{2\pi\sigma^2}}
    \exp
    \left(
        -\frac{(y_i-x_i^\top\beta)^2}{2\sigma^2}
    \right).
\]
Логарифм правдоподобия:
\[
\begin{split}
    \ell(\beta)
    &=
    -\frac{n}{2}\log(2\pi\sigma^2)
    -
    \frac{1}{2\sigma^2}
    \sum_{i=1}^{n}
    (y_i-x_i^\top\beta)^2\\
    &=
    -\frac{n}{2}\log(2\pi\sigma^2)
    -
    \frac{1}{2\sigma^2}
    \|y-X\beta\|_2^2.
\end{split}
\]
Первое слагаемое не зависит от $\beta$, а множитель перед суммой квадратов
отрицателен. Поэтому максимизация $\ell(\beta)$ эквивалентна минимизации
\[
    \|y-X\beta\|_2^2.
\]
\end{proof}

Если $\sigma^2$ также неизвестно и минимальная сумма квадратов остатков
положительна, то
\[
    \widehat\sigma_{\mathrm{ML}}^2
    =
    \frac{\operatorname{RSS}}{n}.
\]
Несмещённая оценка дисперсии использует знаменатель $n-p$:
\[
    s^2=\frac{\operatorname{RSS}}{n-p}.
\]

\subsection{Теорема Гаусса--Маркова}

\begin{definition}
Оценка называется \emph{линейной}, если она имеет вид
\[
    \widetilde\beta=Ay,
\]
где матрица $A$ не зависит от $y$.

Линейная оценка несмещённа для всех $\beta$, если
\[
    AX=I_p.
\]
\end{definition}

\begin{theorem}[Гаусса--Маркова]
Пусть
\[
    y=X\beta+\varepsilon,
\]
где $X\in\R^{n\times p}$ имеет полный столбцовый ранг,
\[
    \E\varepsilon=0,
    \qquad
    \Cov(\varepsilon)=\sigma^2I_n.
\]
Тогда оценка метода наименьших квадратов
\[
    \widehat\beta
    =
    (X^\top X)^{-1}X^\top y
\]
имеет наименьшую ковариационную матрицу среди всех линейных несмещённых
оценок.

Точнее, для любой линейной несмещённой оценки
$\widetilde\beta=Ay$ матрица
\[
    \Cov(\widetilde\beta)
    -
    \Cov(\widehat\beta)
\]
неотрицательно определена. Иными словами, для любого $c\in\R^p$
\[
    \Var(c^\top\widetilde\beta)
    \geq
    \Var(c^\top\widehat\beta).
\]
\end{theorem}

\begin{proof}
Обозначим
\[
    A_0=(X^\top X)^{-1}X^\top,
\]
так что
\[
    \widehat\beta=A_0y.
\]
Поскольку обе оценки несмещённы,
\[
    AX=I_p,
    \qquad
    A_0X=I_p.
\]
Положим
\[
    D=A-A_0.
\]
Тогда
\[
    DX=AX-A_0X=0.
\]
Имеем
\[
    A=A_0+D.
\]
Ковариация произвольной оценки:
\[
\begin{split}
    \Cov(\widetilde\beta)
    &=
    \sigma^2AA^\top\\
    &=
    \sigma^2(A_0+D)(A_0+D)^\top\\
    &=
    \sigma^2
    \left(
        A_0A_0^\top
        +A_0D^\top
        +DA_0^\top
        +DD^\top
    \right).
\end{split}
\]
Покажем, что смешанные слагаемые равны нулю:
\[
\begin{split}
    A_0D^\top
    &=
    (X^\top X)^{-1}X^\top D^\top\\
    &=
    (X^\top X)^{-1}(DX)^\top
    =
    0.
\end{split}
\]
Аналогично $DA_0^\top=0$. Поэтому
\[
    \Cov(\widetilde\beta)
    =
    \sigma^2A_0A_0^\top+\sigma^2DD^\top.
\]
Первое слагаемое равно $\Cov(\widehat\beta)$, следовательно,
\[
    \Cov(\widetilde\beta)
    -
    \Cov(\widehat\beta)
    =
    \sigma^2DD^\top.
\]
Для любого $c\in\R^p$
\[
    c^\top(\sigma^2DD^\top)c
    =
    \sigma^2\|D^\top c\|_2^2
    \geq 0.
\]
Значит, разность ковариационных матриц неотрицательно определена.
\end{proof}

\begin{remark}
Теорема утверждает оптимальность только в классе \emph{линейных
несмещённых} оценок. Она не утверждает, что метод наименьших квадратов
оптимален среди всех возможных оценок без дополнительных предположений.
\end{remark}

\subsection{Распределение оценки и статистический вывод}

\begin{theorem}
Пусть
\[
    y\sim\mathcal{N}(X\beta,\sigma^2I_n),
    \qquad
    \rank X=p<n.
\]
Тогда:

\begin{enumerate}
    \item
    \[
        \widehat\beta
        \sim
        \mathcal{N}
        \left(
            \beta,\,
            \sigma^2(X^\top X)^{-1}
        \right);
    \]
    \item
    \[
        \frac{\operatorname{RSS}}{\sigma^2}
        \sim
        \chi^2_{n-p};
    \]
    \item $\widehat\beta$ и $\operatorname{RSS}$ независимы;
    \item
    \[
        s^2=\frac{\operatorname{RSS}}{n-p}
    \]
    является несмещённой оценкой $\sigma^2$.
\end{enumerate}
\end{theorem}

\begin{proof}
Пусть
\[
    P=X(X^\top X)^{-1}X^\top.
\]
Тогда
\[
    \widehat\beta
    =
    \beta+
    (X^\top X)^{-1}X^\top\varepsilon.
\]
Линейное преобразование нормального вектора нормально, поэтому
\[
    \widehat\beta
    \sim
    \mathcal{N}
    \left(
        \beta,\,
        \sigma^2(X^\top X)^{-1}
    \right).
\]

Остатки равны
\[
    e=y-X\widehat\beta
    =
    (I-P)y
    =
    (I-P)\varepsilon,
\]
так как $(I-P)X=0$. Матрица $I-P$ является ортогональным проектором
ранга $n-p$. Поэтому существует ортогональная матрица $U$, для которой
\[
    I-P
    =
    U
    \begin{pmatrix}
        I_{n-p}&0\\
        0&0_p
    \end{pmatrix}
    U^\top.
\]
Вектор
\[
    z=\frac{1}{\sigma}U^\top\varepsilon
\]
имеет распределение $\mathcal{N}(0,I_n)$. Следовательно,
\[
\begin{split}
    \frac{\operatorname{RSS}}{\sigma^2}
    &=
    \frac{\varepsilon^\top(I-P)\varepsilon}{\sigma^2}\\
    &=
    \sum_{j=1}^{n-p}z_j^2
    \sim
    \chi^2_{n-p}.
\end{split}
\]

Оценка $\widehat\beta-\beta$ зависит от проекции ошибки на
$\col(X)$, а остатки --- от проекции на ортогональное дополнение
$\col(X)^\perp$. Их ковариация равна
\[
\begin{split}
    \Cov(\widehat\beta-\beta,e)
    &=
    \sigma^2
    (X^\top X)^{-1}X^\top(I-P)\\
    &=0.
\end{split}
\]
Обе величины являются линейными преобразованиями совместно нормального
вектора. Для совместно нормальных векторов нулевая ковариация означает
независимость. Следовательно, $\widehat\beta$ и $e$, а значит и
$\operatorname{RSS}=\|e\|^2$, независимы.

Наконец, математическое ожидание случайной величины
$\chi^2_{n-p}$ равно $n-p$. Поэтому
\[
    \E[\operatorname{RSS}]
    =
    (n-p)\sigma^2
\]
и
\[
    \E[s^2]
    =
    \E
    \left[
        \frac{\operatorname{RSS}}{n-p}
    \right]
    =
    \sigma^2.
\]
\end{proof}

Для отдельного коэффициента справедлива статистика
\[
    T_j
    =
    \frac{
        \widehat\beta_j-\beta_j
    }{
        s\sqrt{[(X^\top X)^{-1}]_{jj}}
    }
    \sim
    t_{n-p}.
\]
Отсюда доверительный интервал уровня $1-\alpha$:
\[
    \widehat\beta_j
    \pm
    t_{1-\alpha/2,n-p}
    s\sqrt{[(X^\top X)^{-1}]_{jj}}.
\]

Для среднего ответа в точке $x_0$:
\[
    x_0^\top\widehat\beta
    \pm
    t_{1-\alpha/2,n-p}
    s\sqrt{
        x_0^\top(X^\top X)^{-1}x_0
    }.
\]

Для нового индивидуального наблюдения добавляется неопределённость новой
ошибки:
\[
    x_0^\top\widehat\beta
    \pm
    t_{1-\alpha/2,n-p}
    s\sqrt{
        1+x_0^\top(X^\top X)^{-1}x_0
    }.
\]

\subsection{Регуляризованная линейная регрессия}

\subsubsection{Гребневая регрессия}

Гребневая регрессия, или ridge-регрессия, минимизирует
\[
    \|y-X\beta\|_2^2+\lambda\|\beta\|_2^2,
    \qquad \lambda>0.
\]

\begin{proposition}
Для $\lambda>0$ задача гребневой регрессии имеет единственное решение
\[
    \boxed{
    \widehat\beta_{\mathrm{ridge}}
    =
    (X^\top X+\lambda I)^{-1}X^\top y
    }.
\]
\end{proposition}

\begin{proof}
Градиент целевой функции равен
\[
    2X^\top(X\beta-y)+2\lambda\beta.
\]
Приравнивая его к нулю, получаем
\[
    (X^\top X+\lambda I)\beta=X^\top y.
\]
Для любого ненулевого $v$
\[
    v^\top(X^\top X+\lambda I)v
    =
    \|Xv\|_2^2+\lambda\|v\|_2^2>0.
\]
Следовательно, матрица $X^\top X+\lambda I$ положительно определена и
обратима. Решение единственно и имеет указанный вид.
\end{proof}

На практике свободный член обычно не штрафуют. Тогда вместо $I$ используется
матрица
\[
    L=\diag(0,1,\ldots,1).
\]

Увеличение $\lambda$:

\begin{itemize}
    \item уменьшает абсолютные значения коэффициентов;
    \item повышает устойчивость при мультиколлинеарности;
    \item уменьшает дисперсию оценки;
    \item вносит смещение.
\end{itemize}

\subsubsection{Lasso-регрессия}

Lasso использует штраф по норме $\ell_1$:
\[
    \widehat\beta_{\mathrm{lasso}}
    \in
    \argmin_{\beta}
    \left\{
        \frac{1}{2n}\|y-X\beta\|_2^2
        +
        \lambda\sum_{j=1}^{p}|\beta_j|
    \right\}.
\]
Такой штраф может обращать часть коэффициентов точно в ноль и тем самым
выполнять отбор признаков.

Определим оператор мягкого порога:
\[
    S(z,\lambda)
    =
    \sign(z)(|z|-\lambda)_+.
\]

\begin{algorithmbox}{покоординатный спуск для Lasso}
Предположим, что свободный член не штрафуется.
\begin{enumerate}
    \item Инициализировать коэффициенты $\beta$.
    \item Для каждого признака $j$:
    \begin{enumerate}
        \item вычислить частичный остаток
              \[
                  r_j
                  =
                  y-\beta_0\Ind-\sum_{\ell\ne j}x_\ell\beta_\ell;
              \]
        \item вычислить
              \[
                  z_j=\frac{1}{n}x_j^\top r_j;
              \]
        \item обновить
              \[
                  \beta_j
                  \leftarrow
                  \frac{
                      S(z_j,\lambda)
                  }{
                      \frac{1}{n}\|x_j\|_2^2
                  }.
              \]
    \end{enumerate}
    \item Обновить свободный член как средний остаток.
    \item Повторять проходы до сходимости.
\end{enumerate}
\end{algorithmbox}

\begin{example}
Если признак нормирован так, что
\[
    \frac{1}{n}\|x_j\|^2=1,
\]
и
\[
    z_j=1.4,\qquad \lambda=0.5,
\]
то
\[
    \beta_j
    =
    S(1.4,0.5)
    =
    0.9.
\]
Если бы $|z_j|\leq 0.5$, коэффициент стал бы равен нулю.
\end{example}

\subsection{Метрики качества регрессии}

Среднеквадратичная ошибка:
\[
    \operatorname{MSE}
    =
    \frac{1}{n}
    \sum_{i=1}^{n}
    (y_i-\widehat y_i)^2.
\]

Корень из среднеквадратичной ошибки:
\[
    \operatorname{RMSE}
    =
    \sqrt{\operatorname{MSE}}.
\]

Средняя абсолютная ошибка:
\[
    \operatorname{MAE}
    =
    \frac{1}{n}
    \sum_{i=1}^{n}
    |y_i-\widehat y_i|.
\]

Коэффициент детерминации:
\[
    R^2
    =
    1-
    \frac{
        \sum_{i=1}^{n}(y_i-\widehat y_i)^2
    }{
        \sum_{i=1}^{n}(y_i-\overline y)^2
    }.
\]
При наличии свободного члена и оценивании на обучающей выборке
$0\leq R^2\leq 1$. На тестовой выборке $R^2$ может быть отрицательным.

Скорректированный коэффициент детерминации:
\[
    R_{\mathrm{adj}}^2
    =
    1-
    (1-R^2)
    \frac{n-1}{n-p}.
\]

\subsection{Диагностика и типичные проблемы}

\paragraph{Нелинейность.}
Если зависимость нелинейна по признакам, можно добавить преобразования:
\[
    x^2,\quad x^3,\quad \log x,\quad \sqrt{x},
\]
а также взаимодействия признаков
\[
    x_jx_k.
\]

\paragraph{Мультиколлинеарность.}
Сильная зависимость между столбцами $X$ приводит к большим дисперсиям
коэффициентов и численной неустойчивости. Возможные решения:

\begin{itemize}
    \item удалить избыточные признаки;
    \item объединить коррелированные признаки;
    \item применить снижение размерности;
    \item использовать ridge-регрессию.
\end{itemize}

\paragraph{Гетероскедастичность.}
Если
\[
    \Var(\varepsilon_i\mid x_i)
\]
зависит от объекта, стандартные формулы дисперсий коэффициентов становятся
некорректными. Используются робастные стандартные ошибки, преобразования
целевой переменной или взвешенный метод наименьших квадратов:
\[
    \widehat\beta_{\mathrm{WLS}}
    =
    \argmin_\beta
    (y-X\beta)^\top W(y-X\beta).
\]
При положительно определённой матрице $W$:
\[
    \widehat\beta_{\mathrm{WLS}}
    =
    (X^\top WX)^{-1}X^\top Wy.
\]

\paragraph{Выбросы.}
Квадратичная функция потерь сильно увеличивает влияние наблюдений с большими
остатками. Возможные решения:

\begin{itemize}
    \item проверка качества данных;
    \item робастные функции потерь;
    \item анализ наблюдений с большим рычагом;
    \item преобразование переменных.
\end{itemize}

\paragraph{Экстраполяция.}
Линейная модель может давать неадекватные прогнозы за пределами диапазона
обучающих данных. Хорошее качество интерполяции не гарантирует хорошей
экстраполяции.

% ================================================================
\section{Кластеризация}
% ================================================================

\subsection{Постановка задачи}

Пусть дана выборка без меток:
\[
    X=\{x_1,\ldots,x_n\},
    \qquad
    x_i\in\R^d.
\]

\begin{definition}
\emph{Кластеризация} --- разбиение или вероятностное группирование объектов
так, чтобы объекты внутри одной группы были в некотором смысле похожи,
а объекты из разных групп --- различны.
\end{definition}

При жёсткой кластеризации строится разбиение
\[
    X=C_1\sqcup\cdots\sqcup C_K,
\]
где
\[
    C_j\ne\varnothing,\qquad
    C_i\cap C_j=\varnothing
    \quad\text{при }i\ne j.
\]

Основные виды кластеризации:

\begin{itemize}
    \item \textbf{жёсткая}: каждый объект относится ровно к одному кластеру;
    \item \textbf{мягкая}: объекту сопоставляются вероятности или степени
          принадлежности к нескольким кластерам;
    \item \textbf{иерархическая}: строится система вложенных разбиений;
    \item \textbf{плотностная}: кластеры определяются как области высокой
          плотности, разделённые областями низкой плотности;
    \item \textbf{модельная}: предполагается, что данные порождены смесью
          вероятностных распределений.
\end{itemize}

Кластеризация не имеет единственного универсального определения
``правильного кластера''. Результат зависит от:

\begin{itemize}
    \item выбранных признаков;
    \item масштаба признаков;
    \item функции расстояния;
    \item алгоритма;
    \item гиперпараметров;
    \item предметной интерпретации.
\end{itemize}

\subsection{Меры расстояния и сходства}

Евклидово расстояние:
\[
    d_2(x,z)
    =
    \sqrt{
        \sum_{j=1}^{d}(x_j-z_j)^2
    }.
\]

Манхэттенское расстояние:
\[
    d_1(x,z)
    =
    \sum_{j=1}^{d}|x_j-z_j|.
\]

Расстояние Минковского:
\[
    d_p(x,z)
    =
    \left(
        \sum_{j=1}^{d}|x_j-z_j|^p
    \right)^{1/p}.
\]

Косинусное расстояние:
\[
    d_{\cos}(x,z)
    =
    1-
    \frac{x^\top z}{\|x\|_2\|z\|_2}.
\]

Расстояние Махаланобиса:
\[
    d_M(x,z)
    =
    \sqrt{
        (x-z)^\top\Sigma^{-1}(x-z)
    }.
\]

Если признаки имеют разные единицы измерения, обычно выполняют
стандартизацию:
\[
    x_{ij}^{\mathrm{std}}
    =
    \frac{x_{ij}-\overline x_j}{s_j}.
\]
Без масштабирования признак с большим числовым диапазоном может полностью
определять результат кластеризации.

% ----------------------------------------------------------------
\subsection{Метод $K$-средних}
% ----------------------------------------------------------------

\subsubsection{Целевая функция}

Для каждого объекта введём номер кластера
\[
    c_i\in\{1,\ldots,K\}
\]
и центры
\[
    \mu_1,\ldots,\mu_K\in\R^d.
\]
Метод $K$-средних минимизирует внутрикластерную сумму квадратов:
\[
    J(c,\mu)
    =
    \sum_{i=1}^{n}
    \|x_i-\mu_{c_i}\|_2^2.
\]

Эквивалентная запись:
\[
    J
    =
    \sum_{k=1}^{K}
    \sum_{x_i\in C_k}
    \|x_i-\mu_k\|_2^2.
\]

\begin{lemma}[Оптимальный центр кластера]
Пусть $C=\{x_1,\ldots,x_m\}\subset\R^d$ --- непустой кластер.
Функция
\[
    F(\mu)=\sum_{i=1}^{m}\|x_i-\mu\|_2^2
\]
имеет единственный минимум в точке
\[
    \overline x
    =
    \frac{1}{m}\sum_{i=1}^{m}x_i.
\]
\end{lemma}

\begin{proof}
Для любого $\mu\in\R^d$ запишем
\[
    x_i-\mu
    =
    (x_i-\overline x)+(\overline x-\mu).
\]
Тогда
\[
\begin{split}
    \sum_{i=1}^{m}\|x_i-\mu\|_2^2
    &=
    \sum_{i=1}^{m}
    \left\|
        (x_i-\overline x)+(\overline x-\mu)
    \right\|_2^2\\
    &=
    \sum_{i=1}^{m}\|x_i-\overline x\|_2^2\\
    &\quad+
    2(\overline x-\mu)^\top
    \sum_{i=1}^{m}(x_i-\overline x)\\
    &\quad+
    m\|\overline x-\mu\|_2^2.
\end{split}
\]
Но
\[
    \sum_{i=1}^{m}(x_i-\overline x)
    =
    \sum_{i=1}^{m}x_i-m\overline x
    =
    0.
\]
Следовательно,
\[
    \sum_{i=1}^{m}\|x_i-\mu\|_2^2
    =
    \sum_{i=1}^{m}\|x_i-\overline x\|_2^2
    +
    m\|\overline x-\mu\|_2^2.
\]
Второе слагаемое неотрицательно и равно нулю только при
$\mu=\overline x$. Значит, среднее является единственным минимумом.
\end{proof}

\subsubsection{Алгоритм Ллойда}

\begin{algorithmbox}{метод $K$-средних}
\begin{enumerate}
    \item Задать число кластеров $K$.
    \item Инициализировать центры
          \[
              \mu_1^{(0)},\ldots,\mu_K^{(0)}.
          \]
    \item На итерации $t$ выполнить шаг присваивания:
          \[
              c_i^{(t)}
              =
              \argmin_{1\leq k\leq K}
              \|x_i-\mu_k^{(t)}\|_2^2.
          \]
    \item Выполнить шаг обновления центров:
          \[
              \mu_k^{(t+1)}
              =
              \frac{
                  \sum_{i=1}^{n}
                  \mathbf{1}\{c_i^{(t)}=k\}x_i
              }{
                  \sum_{i=1}^{n}
                  \mathbf{1}\{c_i^{(t)}=k\}
              }.
          \]
    \item Повторять шаги 3--4, пока:
    \begin{itemize}
        \item номера кластеров не перестанут меняться;
        \item либо смещение центров не станет меньше порога;
        \item либо не будет достигнуто максимальное число итераций.
    \end{itemize}
\end{enumerate}
\end{algorithmbox}

Если кластер оказался пустым, его центр нельзя вычислить по формуле среднего.
Возможные практические решения:

\begin{itemize}
    \item перенести центр в наиболее удалённый объект;
    \item повторно инициализировать пустой кластер;
    \item сохранить прежний центр;
    \item перезапустить алгоритм.
\end{itemize}

\subsubsection{Инициализация $K$-means++}

\begin{algorithmbox}{инициализация $K$-means++}
\begin{enumerate}
    \item Выбрать первый центр равновероятно из объектов выборки.
    \item Для каждого объекта $x$ вычислить расстояние до ближайшего уже
          выбранного центра:
          \[
              D(x)
              =
              \min_{\mu\in M}\|x-\mu\|_2,
          \]
          где $M$ --- множество выбранных центров.
    \item Выбрать следующий центр из объектов с вероятностью
          \[
              \PP(x\text{ выбран})
              =
              \frac{D(x)^2}{
                  \sum_{i=1}^{n}D(x_i)^2
              }.
          \]
    \item Повторять шаги 2--3, пока не будет выбрано $K$ центров.
    \item Запустить обычный алгоритм $K$-средних.
\end{enumerate}
\end{algorithmbox}

Инициализация $K$-means++ разносит начальные центры по пространству и обычно
уменьшает чувствительность к неудачному начальному выбору. Для повышения
надёжности алгоритм всё равно запускают несколько раз с разными
инициализациями и выбирают результат с наименьшим значением $J$.

\subsubsection{Монотонность метода $K$-средних}

\begin{theorem}
Каждый шаг алгоритма $K$-средних не увеличивает значение целевой функции
\[
    J(c,\mu)
    =
    \sum_{i=1}^{n}\|x_i-\mu_{c_i}\|_2^2.
\]
Следовательно, последовательность значений $J$ сходится.

Если отсутствуют пустые кластеры и ничьи при выборе ближайшего центра, то
алгоритм заканчивается за конечное число итераций в конфигурации, которая
является покоординатным минимумом по присваиваниям и центрам.
\end{theorem}

\begin{proof}
При фиксированных центрах $\mu_1,\ldots,\mu_K$ алгоритм присваивает
каждому объекту ближайший центр:
\[
    c_i
    =
    \argmin_k\|x_i-\mu_k\|^2.
\]
Поэтому каждый отдельный член суммы принимает наименьшее возможное
значение, и шаг присваивания не увеличивает $J$.

При фиксированном разбиении $C_1,\ldots,C_K$ целевая функция распадается:
\[
    J
    =
    \sum_{k=1}^{K}
    \sum_{x_i\in C_k}
    \|x_i-\mu_k\|^2.
\]
По лемме об оптимальном центре каждое внутреннее выражение минимизируется
при
\[
    \mu_k
    =
    \frac{1}{|C_k|}
    \sum_{x_i\in C_k}x_i.
\]
Значит, шаг обновления центров также не увеличивает $J$.

Таким образом,
\[
    J^{(t+1)}\leq J^{(t)}.
\]
Поскольку $J^{(t)}\geq 0$, монотонная ограниченная снизу последовательность
значений целевой функции сходится.

Возможных присваиваний $n$ объектов к $K$ кластерам не более $K^n$, то есть
их конечное число. При отсутствии ничьих изменение присваивания хотя бы
одного объекта строго уменьшает целевую функцию. Если присваивания не
меняются, но хотя бы один центр не совпадает со средним своего кластера,
обновление центра также строго уменьшает $J$. Поэтому алгоритм не может
бесконечно переходить между различными конфигурациями и заканчивается в
состоянии, в котором:

\begin{itemize}
    \item каждый объект относится к ближайшему центру;
    \item каждый центр является средним объектов своего кластера.
\end{itemize}

Это и есть покоординатный минимум.
\end{proof}

\begin{remark}
Алгоритм не гарантирует нахождение глобального минимума. Разные начальные
центры могут приводить к различным локальным решениям.
\end{remark}

\subsubsection{Минимальный пример}

Пусть
\[
    X=\{1,2,8,9\},
    \qquad K=2,
\]
а начальные центры равны
\[
    \mu_1^{(0)}=1,
    \qquad
    \mu_2^{(0)}=8.
\]

Шаг присваивания:
\[
    C_1=\{1,2\},
    \qquad
    C_2=\{8,9\}.
\]
Обновление центров:
\[
    \mu_1^{(1)}
    =
    \frac{1+2}{2}
    =
    1.5,
\]
\[
    \mu_2^{(1)}
    =
    \frac{8+9}{2}
    =
    8.5.
\]
При следующем присваивании кластеры не меняются, поэтому алгоритм
останавливается.

Начальное значение целевой функции:
\[
    J^{(0)}
    =
    (1-1)^2+(2-1)^2+(8-8)^2+(9-8)^2
    =
    2.
\]
После обновления:
\[
\begin{split}
    J^{(1)}
    &=
    (1-1.5)^2+(2-1.5)^2\\
    &\quad+
    (8-8.5)^2+(9-8.5)^2\\
    &=1.
\end{split}
\]

\begin{figure}[ht]
\centering
\begin{tikzpicture}[xscale=1.1]
    \draw[->] (0,0)--(10,0) node[right] {$x$};

    \foreach \x in {1,2,8,9}
        \fill (\x,0) circle (2.5pt);

    \draw[red,very thick] (1,-0.25)--(1,0.25);
    \draw[blue,very thick] (8,-0.25)--(8,0.25);

    \node[red,above] at (1,0.25) {$\mu_1^{(0)}$};
    \node[blue,above] at (8,0.25) {$\mu_2^{(0)}$};

    \draw[red,very thick,dashed] (1.5,-0.65)--(1.5,-0.25);
    \draw[blue,very thick,dashed] (8.5,-0.65)--(8.5,-0.25);

    \node[red,below] at (1.5,-0.65) {$\mu_1^{(1)}=1.5$};
    \node[blue,below] at (8.5,-0.65) {$\mu_2^{(1)}=8.5$};
\end{tikzpicture}
\caption{Одна итерация метода $K$-средних.}
\end{figure}

\subsubsection{Свойства и ограничения}

Преимущества:

\begin{itemize}
    \item простота;
    \item высокая скорость;
    \item понятные центры кластеров;
    \item хорошая масштабируемость.
\end{itemize}

Ограничения:

\begin{itemize}
    \item требуется заранее задать $K$;
    \item чувствительность к масштабу признаков;
    \item чувствительность к выбросам;
    \item зависимость от инициализации;
    \item предпочтение компактных, приблизительно сферических кластеров;
    \item использование только евклидовой геометрии и средних значений.
\end{itemize}

Одна итерация требует порядка
\[
    O(nKd)
\]
операций. При $T$ итерациях сложность составляет
\[
    O(nKdT).
\]

% ----------------------------------------------------------------
\subsection{Выбор числа кластеров}
% ----------------------------------------------------------------

\subsubsection{Метод локтя}

Для разных значений $K$ вычисляют
\[
    J_K
    =
    \sum_{k=1}^{K}
    \sum_{x_i\in C_k}
    \|x_i-\mu_k\|^2.
\]
Функция $J_K$ не возрастает с увеличением $K$. Ищут точку, после которой
дальнейшее увеличение числа кластеров даёт относительно небольшое уменьшение
$J_K$.

Метод локтя является эвристическим: отчётливый изгиб может отсутствовать.

\subsubsection{Силуэт}

Для объекта $x_i$ определим:

\[
    a(i)
    =
    \text{среднее расстояние от }x_i
    \text{ до объектов его кластера},
\]
\[
    b(i)
    =
    \min_{C\ne C(x_i)}
    \left\{
        \text{среднее расстояние от }x_i
        \text{ до объектов }C
    \right\}.
\]

Коэффициент силуэта:
\[
    s(i)
    =
    \frac{b(i)-a(i)}{\max\{a(i),b(i)\}}.
\]
При этом
\[
    -1\leq s(i)\leq 1.
\]

Интерпретация:

\begin{itemize}
    \item $s(i)\approx 1$: объект хорошо согласуется со своим кластером;
    \item $s(i)\approx 0$: объект находится на границе кластеров;
    \item $s(i)<0$: объект, возможно, отнесён к неверному кластеру.
\end{itemize}

Средний силуэт:
\[
    \overline s
    =
    \frac{1}{n}\sum_{i=1}^{n}s(i)
\]
может использоваться для сравнения разных значений $K$.

\begin{example}
Пусть кластеры равны
\[
    C_1=\{0,1\},
    \qquad
    C_2=\{5,6\}.
\]
Для точки $0$
\[
    a(0)=1,
\]
\[
    b(0)=\frac{|0-5|+|0-6|}{2}=5.5.
\]
Следовательно,
\[
    s(0)
    =
    \frac{5.5-1}{5.5}
    \approx 0.818.
\]
\end{example}

% ----------------------------------------------------------------
\subsection{Иерархическая кластеризация}
% ----------------------------------------------------------------

Иерархическая кластеризация строит дерево вложенных кластеров ---
\emph{дендрограмму}.

Основные подходы:

\begin{itemize}
    \item \textbf{агломеративный}: начать с отдельных объектов и
          последовательно объединять кластеры;
    \item \textbf{дивизивный}: начать с одного общего кластера и
          последовательно разделять его.
\end{itemize}

\begin{algorithmbox}{агломеративная иерархическая кластеризация}
\begin{enumerate}
    \item Поместить каждый объект в отдельный кластер:
          \[
              C_i=\{x_i\}.
          \]
    \item Вычислить расстояния между всеми парами кластеров.
    \item Найти два ближайших кластера.
    \item Объединить их.
    \item Пересчитать расстояния между новым кластером и остальными.
    \item Повторять шаги 3--5, пока не останется один кластер.
    \item Выбрать уровень разреза дендрограммы и получить требуемое число
          кластеров.
\end{enumerate}
\end{algorithmbox}

Расстояние между кластерами определяется правилом связи.

\paragraph{Одиночная связь:}
\[
    d_{\mathrm{single}}(A,B)
    =
    \min_{x\in A,\;z\in B}d(x,z).
\]

\paragraph{Полная связь:}
\[
    d_{\mathrm{complete}}(A,B)
    =
    \max_{x\in A,\;z\in B}d(x,z).
\]

\paragraph{Средняя связь:}
\[
    d_{\mathrm{average}}(A,B)
    =
    \frac{1}{|A||B|}
    \sum_{x\in A}
    \sum_{z\in B}
    d(x,z).
\]

\paragraph{Метод Уорда.}
Объединяется пара кластеров, вызывающая наименьший прирост внутрикластерной
суммы квадратов:
\[
    \Delta(A,B)
    =
    \frac{|A||B|}{|A|+|B|}
    \|\mu_A-\mu_B\|_2^2.
\]

Одиночная связь способна находить вытянутые кластеры, но подвержена эффекту
``цепочки''. Полная связь формирует более компактные кластеры. Метод Уорда
по поведению близок к минимизации критерия $K$-средних.

\begin{example}
Пусть даны точки
\[
    X=\{0,1,5\}.
\]
При одиночной связи сначала объединяются точки $0$ и $1$, поскольку
\[
    d(0,1)=1.
\]
Расстояние между кластером $\{0,1\}$ и точкой $5$:
\[
    d_{\mathrm{single}}(\{0,1\},\{5\})
    =
    \min\{5,4\}
    =
    4.
\]
Поэтому второе объединение происходит на высоте $4$.

\begin{center}
\begin{tikzpicture}[xscale=1.2,yscale=0.7]
    \draw (0,0)--(0,1);
    \draw (1,0)--(1,1);
    \draw (0,1)--(1,1);
    \draw (0.5,1)--(0.5,4);

    \draw (3,0)--(3,4);
    \draw (0.5,4)--(3,4);

    \node[below] at (0,0) {$0$};
    \node[below] at (1,0) {$1$};
    \node[below] at (3,0) {$5$};

    \draw[->] (-0.5,0)--(-0.5,4.6);
    \node[above] at (-0.5,4.6) {расстояние};

    \node[left] at (-0.5,1) {$1$};
    \node[left] at (-0.5,4) {$4$};
\end{tikzpicture}
\end{center}
\end{example}

Для хранения полной матрицы попарных расстояний требуется
$O(n^2)$ памяти. Наивная реализация может требовать $O(n^3)$ времени;
оптимизированные реализации для распространённых правил связи обычно
имеют квадратичную или близкую к квадратичной сложность.

% ----------------------------------------------------------------
\subsection{DBSCAN}
% ----------------------------------------------------------------

DBSCAN определяет кластеры как плотные области, разделённые разреженными
областями. Алгоритм использует два параметра:

\begin{itemize}
    \item $\varepsilon>0$ --- радиус окрестности;
    \item $\mathrm{MinPts}$ --- минимальное число объектов в окрестности.
\end{itemize}

\begin{definition}
$\varepsilon$-окрестность точки $p$:
\[
    N_\varepsilon(p)
    =
    \{q\in X:d(p,q)\leq\varepsilon\}.
\]
Сама точка $p$ включается в свою окрестность.
\end{definition}

\begin{definition}
Точка $p$ называется \emph{ядерной}, если
\[
    |N_\varepsilon(p)|
    \geq
    \mathrm{MinPts}.
\]
\end{definition}

\begin{definition}
Точка $q$ непосредственно достижима по плотности из $p$, если
\[
    q\in N_\varepsilon(p)
\]
и $p$ является ядерной точкой.
\end{definition}

\begin{definition}
Точка $q$ достижима по плотности из $p$, если существует последовательность
\[
    p=p_1,p_2,\ldots,p_m=q,
\]
в которой $p_{i+1}$ непосредственно достижима по плотности из $p_i$.
\end{definition}

\begin{definition}
Точки $p$ и $q$ связаны по плотности, если существует точка $o$, из которой
и $p$, и $q$ достижимы по плотности.
\end{definition}

Различают:

\begin{itemize}
    \item \textbf{ядерные точки};
    \item \textbf{граничные точки}, лежащие в окрестности ядерной точки,
          но сами не являющиеся ядерными;
    \item \textbf{шум}, не принадлежащий ни одному кластеру.
\end{itemize}

\begin{algorithmbox}{DBSCAN}
\begin{enumerate}
    \item Пометить все объекты как непосещённые.
    \item Выбрать непосещённую точку $p$ и пометить её как посещённую.
    \item Найти $N_\varepsilon(p)$.
    \item Если
          \[
              |N_\varepsilon(p)|<\mathrm{MinPts},
          \]
          временно пометить $p$ как шум.
    \item Иначе:
    \begin{enumerate}
        \item создать новый кластер $C$;
        \item добавить $p$ и точки из $N_\varepsilon(p)$ в множество
              кандидатов;
        \item для каждой точки-кандидата $q$:
        \begin{enumerate}
            \item если $q$ ещё не посещена, пометить её посещённой и найти
                  $N_\varepsilon(q)$;
            \item если $q$ является ядерной, добавить её соседей
                  в множество кандидатов;
            \item если $q$ ещё не принадлежит кластеру, добавить её в $C$.
        \end{enumerate}
    \end{enumerate}
    \item Повторять, пока все точки не будут посещены.
\end{enumerate}
\end{algorithmbox}

Точка, первоначально помеченная как шум, позднее может стать граничной
точкой найденного кластера.

\begin{example}
Пусть
\[
    X=\{0,\,0.1,\,0.2,\,3\},
\]
\[
    \varepsilon=0.15,
    \qquad
    \mathrm{MinPts}=2.
\]
Тогда
\[
    N_\varepsilon(0)=\{0,0.1\},
\]
\[
    N_\varepsilon(0.1)=\{0,0.1,0.2\},
\]
\[
    N_\varepsilon(0.2)=\{0.1,0.2\}.
\]
Первые три точки образуют один кластер. Для точки $3$
\[
    N_\varepsilon(3)=\{3\},
\]
поэтому она является шумом.

\begin{center}
\begin{tikzpicture}[xscale=2]
    \draw[->] (-0.2,0)--(3.6,0) node[right] {$x$};

    \fill[blue] (0,0) circle (3pt);
    \fill[blue] (0.1,0) circle (3pt);
    \fill[blue] (0.2,0) circle (3pt);
    \fill[red] (3,0) circle (3pt);

    \node[above,blue] at (0.1,0.1) {кластер};
    \node[above,red] at (3,0.1) {шум};
\end{tikzpicture}
\end{center}
\end{example}

Преимущества DBSCAN:

\begin{itemize}
    \item не требует заранее задавать число кластеров;
    \item находит кластеры произвольной геометрической формы;
    \item явно выделяет шум;
    \item устойчивее $K$-средних к отдельным выбросам.
\end{itemize}

Ограничения:

\begin{itemize}
    \item чувствительность к выбору $\varepsilon$ и $\mathrm{MinPts}$;
    \item трудности при сильно различающейся плотности кластеров;
    \item ухудшение качества расстояний в пространствах высокой размерности;
    \item зависимость результата от масштабирования признаков.
\end{itemize}

При наличии подходящего пространственного индекса в пространствах небольшой
размерности типичная сложность близка к
\[
    O(n\log n).
\]
В худшем случае она может достигать
\[
    O(n^2).
\]

% ----------------------------------------------------------------
\subsection{Гауссовские смеси и EM-алгоритм}
% ----------------------------------------------------------------

\subsubsection{Вероятностная модель}

В модели смеси гауссовских распределений предполагается
\[
    p(x)
    =
    \sum_{k=1}^{K}
    \pi_k
    \mathcal{N}(x\mid\mu_k,\Sigma_k),
\]
где
\[
    \pi_k\geq 0,
    \qquad
    \sum_{k=1}^{K}\pi_k=1.
\]

Вводится скрытая переменная
\[
    z_i\in\{1,\ldots,K\},
\]
обозначающая компоненту, породившую объект $x_i$:
\[
    \PP(z_i=k)=\pi_k,
\]
\[
    x_i\mid z_i=k
    \sim
    \mathcal{N}(\mu_k,\Sigma_k).
\]

В отличие от $K$-средних объект не обязан жёстко принадлежать одной
компоненте. Вычисляются вероятности принадлежности:
\[
    \gamma_{ik}
    =
    \PP(z_i=k\mid x_i).
\]

\subsubsection{EM-алгоритм}

По формуле Байеса
\[
    \gamma_{ik}
    =
    \frac{
        \pi_k
        \mathcal{N}(x_i\mid\mu_k,\Sigma_k)
    }{
        \sum_{\ell=1}^{K}
        \pi_\ell
        \mathcal{N}(x_i\mid\mu_\ell,\Sigma_\ell)
    }.
\]

Обозначим эффективное число объектов компоненты:
\[
    N_k=\sum_{i=1}^{n}\gamma_{ik}.
\]

\begin{algorithmbox}{EM для смеси гауссовских распределений}
\begin{enumerate}
    \item Инициализировать параметры
          \[
              \pi_k^{(0)},\quad
              \mu_k^{(0)},\quad
              \Sigma_k^{(0)}.
          \]
    \item \textbf{E-шаг.} Вычислить ответственности:
          \[
              \gamma_{ik}
              =
              \frac{
                  \pi_k
                  \mathcal{N}(x_i\mid\mu_k,\Sigma_k)
              }{
                  \sum_{\ell=1}^{K}
                  \pi_\ell
                  \mathcal{N}(x_i\mid\mu_\ell,\Sigma_\ell)
              }.
          \]
    \item \textbf{M-шаг.} Вычислить
          \[
              N_k=\sum_{i=1}^{n}\gamma_{ik}
          \]
          и обновить параметры:
          \[
              \pi_k
              \leftarrow
              \frac{N_k}{n},
          \]
          \[
              \mu_k
              \leftarrow
              \frac{1}{N_k}
              \sum_{i=1}^{n}\gamma_{ik}x_i,
          \]
          \[
              \Sigma_k
              \leftarrow
              \frac{1}{N_k}
              \sum_{i=1}^{n}
              \gamma_{ik}
              (x_i-\mu_k)(x_i-\mu_k)^\top.
          \]
    \item Вычислить логарифм правдоподобия
          \[
              \ell(\theta)
              =
              \sum_{i=1}^{n}
              \log
              \left(
                  \sum_{k=1}^{K}
                  \pi_k
                  \mathcal{N}(x_i\mid\mu_k,\Sigma_k)
              \right).
          \]
    \item Повторять E- и M-шаги, пока изменение логарифма правдоподобия
          не станет меньше заданного порога.
\end{enumerate}
\end{algorithmbox}

\subsubsection{Монотонность EM-алгоритма}

\begin{theorem}
Каждая итерация EM-алгоритма не уменьшает логарифм правдоподобия наблюдаемых
данных:
\[
    \ell(\theta^{(t+1)})
    \geq
    \ell(\theta^{(t)}).
\]
\end{theorem}

\begin{proof}
Для каждого наблюдения $x_i$ введём произвольное распределение
$q_i(z)$ над скрытой переменной. Определим функционал
\[
    \mathcal{F}(q,\theta)
    =
    \sum_{i=1}^{n}
    \sum_z
    q_i(z)
    \log
    \frac{
        p(x_i,z\mid\theta)
    }{
        q_i(z)
    }.
\]

Для каждого $i$ справедливо тождество
\[
    \log p(x_i\mid\theta)
    =
    \sum_zq_i(z)
    \log
    \frac{
        p(x_i,z\mid\theta)
    }{
        q_i(z)
    }
    +
    \operatorname{KL}
    \left(
        q_i(z)
        \,\|\,
        p(z\mid x_i,\theta)
    \right).
\]
Действительно,
\[
\begin{split}
    &\operatorname{KL}
    \left(
        q_i(z)
        \,\|\,
        p(z\mid x_i,\theta)
    \right)\\
    &=
    \sum_zq_i(z)
    \log
    \frac{
        q_i(z)
    }{
        p(z\mid x_i,\theta)
    }\\
    &=
    \sum_zq_i(z)
    \log
    \frac{
        q_i(z)p(x_i\mid\theta)
    }{
        p(x_i,z\mid\theta)
    }\\
    &=
    \log p(x_i\mid\theta)
    -
    \sum_zq_i(z)
    \log
    \frac{
        p(x_i,z\mid\theta)
    }{
        q_i(z)
    }.
\end{split}
\]
Суммируя по $i$, получаем
\[
    \ell(\theta)
    =
    \mathcal{F}(q,\theta)
    +
    \sum_{i=1}^{n}
    \operatorname{KL}
    \left(
        q_i
        \,\|\,
        p(z_i\mid x_i,\theta)
    \right).
\]
Поскольку дивергенция Кульбака--Лейблера неотрицательна,
\[
    \ell(\theta)\geq\mathcal{F}(q,\theta).
\]

На E-шаге выбирается
\[
    q_i^{(t)}(z)
    =
    p(z\mid x_i,\theta^{(t)}).
\]
Тогда все KL-дивергенции равны нулю:
\[
    \ell(\theta^{(t)})
    =
    \mathcal{F}(q^{(t)},\theta^{(t)}).
\]

На M-шаге выбирается параметр, максимизирующий функционал при фиксированном
$q^{(t)}$:
\[
    \mathcal{F}(q^{(t)},\theta^{(t+1)})
    \geq
    \mathcal{F}(q^{(t)},\theta^{(t)}).
\]
Кроме того,
\[
    \ell(\theta^{(t+1)})
    \geq
    \mathcal{F}(q^{(t)},\theta^{(t+1)}).
\]
Объединяя неравенства, получаем
\[
\begin{split}
    \ell(\theta^{(t+1)})
    &\geq
    \mathcal{F}(q^{(t)},\theta^{(t+1)})\\
    &\geq
    \mathcal{F}(q^{(t)},\theta^{(t)})\\
    &=
    \ell(\theta^{(t)}).
\end{split}
\]
\end{proof}

\begin{remark}
Монотонность не означает достижения глобального максимума. EM может
сходиться к локальному максимуму или стационарной точке и зависит от
инициализации.
\end{remark}

\subsubsection{Минимальный пример EM}

Пусть
\[
    X=\{-2,-1,1,2\},
    \qquad K=2.
\]
Начальные параметры:
\[
    \mu_1=-1.5,\qquad
    \mu_2=1.5,
\]
\[
    \sigma_1^2=\sigma_2^2=1,
    \qquad
    \pi_1=\pi_2=\frac12.
\]

На E-шаге ответственности первой компоненты приблизительно равны
\[
\begin{array}{c|cccc}
x_i&-2&-1&1&2\\
\hline
\gamma_{i1}
&0.9975&0.9526&0.0474&0.0025
\end{array}
\]
и
\[
    N_1
    =
    \sum_i\gamma_{i1}
    \approx 2.
\]
Новое среднее первой компоненты:
\[
\begin{split}
    \mu_1^{\mathrm{new}}
    &=
    \frac{
        -2\cdot0.9975
        -1\cdot0.9526
        +1\cdot0.0474
        +2\cdot0.0025
    }{2}\\
    &\approx -1.448.
\end{split}
\]
По симметрии
\[
    \mu_2^{\mathrm{new}}\approx 1.448.
\]
Новые веса остаются близкими к
\[
    \pi_1=\pi_2=0.5,
\]
а обновлённые дисперсии приблизительно равны $0.404$.

\subsubsection{Связь с методом $K$-средних}

Если все компоненты имеют одинаковые сферические ковариации
\[
    \Sigma_k=\sigma^2I,
\]
а распределение принадлежности заменяется жёстким выбором наиболее
вероятной компоненты, то EM становится близок к методу $K$-средних.
При $\sigma^2\to 0$ ответственности стремятся к значениям $0$ или $1$,
и каждый объект относится к ближайшему среднему.

\subsubsection{Ограничения гауссовских смесей}

\begin{itemize}
    \item требуется задать число компонент;
    \item результат зависит от инициализации;
    \item одна гауссовская компонента описывает эллипсоидальную область;
    \item вычисление полных ковариационных матриц дорого при большом $d$;
    \item правдоподобие может становиться неограниченным, если ковариация
          компоненты схлопывается около отдельного объекта.
\end{itemize}

Для предотвращения вырождения используют регуляризацию:
\[
    \Sigma_k
    \leftarrow
    \Sigma_k+\delta I,
    \qquad
    \delta>0.
\]

% ----------------------------------------------------------------
\subsection{Сравнение методов кластеризации}
% ----------------------------------------------------------------

\begin{table}[ht]
\centering
\small
\begin{tabularx}{\textwidth}{
    >{\raggedright\arraybackslash}p{2.6cm}
    >{\raggedright\arraybackslash}p{3.1cm}
    >{\raggedright\arraybackslash}X
    >{\raggedright\arraybackslash}X}
\toprule
\textbf{Метод} &
\textbf{Основная идея} &
\textbf{Преимущества} &
\textbf{Ограничения}\\
\midrule

$K$-средних &
Минимизация суммы квадратов расстояний до центров. &
Быстрота, простота, масштабируемость, понятные центры. &
Нужно задать $K$; чувствительность к масштабу, выбросам и инициализации;
плохо описывает несферические кластеры.\\

Иерархическая кластеризация &
Последовательное объединение или разделение кластеров. &
Дендрограмма; не обязательно сразу фиксировать число кластеров; разные
правила связи. &
Квадратичная память; высокая стоимость на больших выборках; необратимость
ранних объединений.\\

DBSCAN &
Поиск связных областей высокой плотности. &
Кластеры произвольной формы; выделение шума; не требуется число кластеров. &
Сложный выбор параметров; проблемы при различной плотности и высокой
размерности.\\

Гауссовская смесь &
Представление плотности как смеси нормальных распределений. &
Мягкая принадлежность; вероятностная интерпретация; эллиптические кластеры. &
Требуется число компонент; локальные решения; возможное вырождение
ковариаций.\\

\bottomrule
\end{tabularx}
\end{table}

\subsection{Оценивание качества кластеризации}

\subsubsection{Внутренние критерии}

Используют только объекты и найденные кластеры:

\begin{itemize}
    \item внутрикластерную сумму квадратов;
    \item коэффициент силуэта;
    \item отношение внутрикластерных и межкластерных расстояний;
    \item устойчивость результата при изменении выборки и инициализации.
\end{itemize}

Внутренние критерии оценивают геометрические свойства, но не гарантируют
предметную полезность кластеров.

\subsubsection{Внешние критерии}

Если известна эталонная разметка, используют:

\begin{itemize}
    \item индекс Рэнда и скорректированный индекс Рэнда;
    \item нормированную взаимную информацию;
    \item однородность, полноту и V-меру.
\end{itemize}

Номера кластеров сами по себе не имеют смысла: перестановка их номеров не
должна менять оценку качества.

\subsubsection{Предметная проверка}

Найденные группы должны:

\begin{itemize}
    \item иметь содержательную интерпретацию;
    \item быть устойчивыми к разумным изменениям данных;
    \item воспроизводиться на новых выборках;
    \item решать исходную прикладную задачу.
\end{itemize}

Высокое значение математического критерия не означает автоматически, что
кластеры полезны.

\subsection{Практическая схема кластеризации}

\begin{algorithmbox}{проведение кластерного анализа}
\begin{enumerate}
    \item Сформулировать смысл понятия ``похожий объект''.
    \item Выбрать признаки, относящиеся к поставленной задаче.
    \item Обработать пропуски и выбросы.
    \item Масштабировать признаки.
    \item Выбрать расстояние или вероятностную модель.
    \item Запустить несколько подходящих методов и наборов гиперпараметров.
    \item Оценить:
    \begin{itemize}
        \item внутренние критерии;
        \item устойчивость;
        \item согласованность с предметной областью.
    \end{itemize}
    \item Интерпретировать каждый кластер.
    \item Зафиксировать масштабирование, параметры, случайные инициализации
          и правила обработки новых объектов.
\end{enumerate}
\end{algorithmbox}

\subsection{Типичные ошибки}

\begin{enumerate}
    \item Применение евклидова расстояния к признакам разных масштабов.
    \item Кластеризация по идентификаторам или техническим полям.
    \item Интерпретация случайных геометрических групп как реальных классов.
    \item Выбор числа кластеров только по одному критерию.
    \item Единственный запуск алгоритма $K$-средних или EM.
    \item Игнорирование выбросов.
    \item Использование слишком большого количества неинформативных
          признаков.
    \item Сравнение номеров кластеров без учёта возможности их перестановки.
    \item Использование кластеризации как доказательства причинно-следственных
          связей.
\end{enumerate}

% ================================================================
\section{Краткое сопоставление регрессии и кластеризации}
% ================================================================

\begin{table}[ht]
\centering
\begin{tabularx}{\textwidth}{
    >{\raggedright\arraybackslash}p{4cm}
    >{\raggedright\arraybackslash}X
    >{\raggedright\arraybackslash}X}
\toprule
\textbf{Характеристика} &
\textbf{Линейная регрессия} &
\textbf{Кластеризация}\\
\midrule

Тип обучения &
С учителем. &
Без учителя.\\

Данные &
Пары $(x_i,y_i)$. &
Только объекты $x_i$.\\

Цель &
Предсказать числовую переменную. &
Обнаружить структуру и группы объектов.\\

Типичный критерий &
Сумма квадратов ошибок прогноза. &
Компактность, разделённость, плотность или правдоподобие.\\

Проверка &
На объектах с известным истинным ответом. &
По внутренним, внешним и предметным критериям.\\

Интерпретация &
Коэффициенты характеризуют изменение среднего прогноза. &
Кластеры требуют дополнительной предметной интерпретации.\\

Основные риски &
Нелинейность, мультиколлинеарность, выбросы, гетероскедастичность,
переобучение. &
Неподходящее расстояние, масштаб признаков, нестабильность, отсутствие
объективно единственного разбиения.\\

\bottomrule
\end{tabularx}
\end{table}

% ================================================================
\section{Итоговые формулы}
% ================================================================

\subsection*{Линейная регрессия}

\[
    \widehat y=X\widehat\beta,
\]
\[
    \widehat\beta
    =
    (X^\top X)^{-1}X^\top y,
\]
\[
    e=y-X\widehat\beta,
\]
\[
    X^\top e=0,
\]
\[
    \Cov(\widehat\beta\mid X)
    =
    \sigma^2(X^\top X)^{-1},
\]
\[
    s^2
    =
    \frac{\|e\|_2^2}{n-p}.
\]

Ridge:
\[
    \widehat\beta_{\mathrm{ridge}}
    =
    (X^\top X+\lambda I)^{-1}X^\top y.
\]

Градиентный спуск:
\[
    \beta^{(t+1)}
    =
    \beta^{(t)}
    -
    \frac{\eta}{n}
    X^\top
    \bigl(X\beta^{(t)}-y\bigr).
\]

\subsection*{Метод $K$-средних}

\[
    c_i
    =
    \argmin_k\|x_i-\mu_k\|_2^2,
\]
\[
    \mu_k
    =
    \frac{
        \sum_i\mathbf{1}\{c_i=k\}x_i
    }{
        \sum_i\mathbf{1}\{c_i=k\}
    },
\]
\[
    J
    =
    \sum_i\|x_i-\mu_{c_i}\|_2^2.
\]

\subsection*{Силуэт}

\[
    s(i)
    =
    \frac{b(i)-a(i)}{\max\{a(i),b(i)\}}.
\]

\subsection*{Гауссовская смесь}

\[
    p(x)
    =
    \sum_{k=1}^{K}
    \pi_k\mathcal{N}(x\mid\mu_k,\Sigma_k),
\]
\[
    \gamma_{ik}
    =
    \frac{
        \pi_k\mathcal{N}(x_i\mid\mu_k,\Sigma_k)
    }{
        \sum_{\ell}
        \pi_\ell\mathcal{N}(x_i\mid\mu_\ell,\Sigma_\ell)
    },
\]
\[
    N_k=\sum_i\gamma_{ik},
\]
\[
    \pi_k=\frac{N_k}{n},
\qquad
    \mu_k=\frac{1}{N_k}\sum_i\gamma_{ik}x_i,
\]
\[
    \Sigma_k
    =
    \frac{1}{N_k}
    \sum_i
    \gamma_{ik}
    (x_i-\mu_k)(x_i-\mu_k)^\top.
\]

\end{document}
